\documentclass[11pt]{article}
\usepackage{amsmath,amssymb}
\usepackage[margin=2.3cm]{geometry}

\title{Cálculos tipo «hechos a mano» con \texttt{handcalcs} + \texttt{pint}}
\author{}
\date{\today}

\begin{document}
\maketitle

Con \texttt{handcalcs} cada cálculo se muestra como en un examen: la fórmula, la
sustitución de valores y el resultado. El patrón en Calc es:

\begin{enumerate}
  \item En una celda defines una función de cálculo y la pasas a \verb|hc(func, args)|.
  \item \verb|hc| devuelve el LaTeX (no se imprime: solo aparece en el documento).
  \item En el cuerpo escribes \verb|\py{nombre}| donde quieras el resultado.
\end{enumerate}

\section{Viga simplemente apoyada}
Momento y cortante máximos de una viga con carga uniforme:

%#python
def viga(w, L):
    M_max = w * L**2 / 8      # momento máximo
    V_max = w * L / 2         # cortante máximo
    return M_max, V_max

# hc(...) devuelve LaTeX listo (display math); no se imprime -> sin doble render
calc_viga = hc(viga, 12.0, 6.0)
%#end

\py{calc_viga}

con $w = 12\ \mathrm{kN/m}$ y $L = 6\ \mathrm{m}$.

\section{Cálculo con unidades (\texttt{pint})}
\texttt{pint} arrastra las unidades automáticamente y \texttt{handcalcs} las
muestra en el resultado:

%#python
def energia(m, v):
    E_c = (m * v**2) / 2       # energía cinética
    return E_c

calc_energia = hc(energia, 1500 * ureg.kg, 25 * ureg.m / ureg.s)
%#end

\py{calc_energia}

\section{Tensión normal}
Combinando varias líneas en un solo bloque renderizado:

%#python
def tension(P, b, h):
    A = b * h          # área de la sección rectangular
    sigma = P / A      # tensión normal
    return sigma

calc_tension = hc(tension, 50 * ureg.kN, 20 * ureg.mm, 30 * ureg.mm)
%#end

\py{calc_tension}

\medskip
Todo esto se calcula en Python y se tipografía en \LaTeX{} sin salir del documento.

\end{document}
